\documentclass[reqno]{amsart}
\usepackage{amsmath,amssymb,amsthm,mathtools}
\usepackage[margin=0.92in]{geometry}
\usepackage{enumitem}
\setlist{nosep,leftmargin=*}
\usepackage{booktabs}
\usepackage[colorlinks=true,linkcolor=blue,citecolor=blue]{hyperref}

\theoremstyle{plain}
\newtheorem{theorem}{Theorem}[section]
\newtheorem{proposition}[theorem]{Proposition}
\newtheorem{lemma}[theorem]{Lemma}
\newtheorem{corollary}[theorem]{Corollary}
\theoremstyle{definition}
\newtheorem{definition}[theorem]{Definition}
\newtheorem{example}[theorem]{Example}
\newtheorem{remark}[theorem]{Remark}
\newtheorem{convention}[theorem]{Convention}

\DeclareMathOperator{\soc}{soc}
\DeclareMathOperator{\topp}{top}
\DeclareMathOperator{\Hom}{Hom}
\DeclareMathOperator{\Ext}{Ext}
\DeclareMathOperator{\End}{End}
\newcommand{\dimv}{\underline{\dim}\,}
\newcommand{\la}{\lambda}
\newcommand{\al}{\alpha}
\newcommand{\om}{\omega}
\newcommand{\ka}{\kappa}
\newcommand{\Pia}{\Pi}

\makeatletter
\def\@secnumfont{\bfseries}
\makeatother
\setlength{\topsep}{2pt plus 1pt minus 1pt}

\begin{document}

\title[Type $E$ examples with $c_p=3$]{Type $E$ examples beyond the
$\la$-minuscule range:\\ the extremal weights with a $c$-trace entry $3$}

\author{Jaideep Sai Padhi}
\address{Department of Mathematics, Purdue University, West Lafayette, IN 47907}
\email{jpadhi@purdue.edu}

\begin{abstract}
Let $\Gamma$ be a simply-laced Dynkin diagram, $\la=\om_\ka$ a fundamental
weight, $w$ an element of the Weyl group of minimal length in $wW_\la$, and
$V$ the unique $T$-fixed representation on
$\mathfrak{M}(\la-w\la,\la)$.  The \emph{$c$-trace} of $w$ is the sequence
$c_p=\langle\al_{i_p}^\vee,\mu^{(p-1)}\rangle$ attached to a reduced word, and
$w$ is $\la$-minuscule exactly when every $c_p=1$.  We determine, by an
exhaustive computation over all of $E_6$ and $E_7$ and all fundamental
$\la$, every extremal weight whose $c$-trace attains the value $3$ and whose
$V$ has simple spectrum.  There are exactly nine, listed in Table
\ref{tab:nine}; one of them, number (4), is not fully commutative.  No
extremal weight in $E_6\cup E_7$ with a $c$-trace entry $\ge4$ has simple
spectrum.  The computation is carried out twice, by independent methods, and
the claims of this note are computational, not proved.
\end{abstract}

\maketitle

\section{What is computed, and what is not}\label{sec:scope}

Everything in this note is the output of a finite computation, performed in two
independent ways that agree on every case where both terminate.  Nothing here
is proved.  We say this at the outset because the companion
work is proof-based and we do not wish the two to be confused.

The two methods are:

\begin{enumerate}[label=(\roman*)]
\item \emph{Socle filtration.}  Build the injective $\Pia$-module $I_\ka$
explicitly as $\Pia e_{\nu(\ka)}$, where $\nu$ is the Nakayama permutation,
and compute the iterated socle $\soc_\sigma(I_\ka)$ along the $c$-weighted
sequence $\sigma$ of a reduced word for $w$.  Each socle step of colour $j$ is
the kernel computation
\[
   X\ \longmapsto\ \{\,x\in I_\ka:\ a_{j\to i}\,x\in X\ \text{for every } i\sim j\,\},
\]
so the whole filtration costs one linear-algebra pass per letter.  By the main
theorem of the companion work, $V\cong\soc_\sigma(I_\ka)$, so the graded
dimensions of the socle layers are the graded dimensions of $V$, and simple
spectrum is the condition that every layer be multiplicity-free in each colour.
\item \emph{Submodule search.}  Search directly for a graded submodule of
$I_\ka$ of dimension vector $v=\la-w\la$, by depth-first search over the
subspace lattice degree by degree.  This uses no theorem; it is slow, and in
$E_7$ it does not terminate on the larger weights.
\end{enumerate}

We validated (i) against (ii) on $1001$ extremal weights --- all of
$A_3\,\om_1$, $D_4\,\om_2$, $D_5\,\om_4$, $E_6\,\om_4$, $E_6\,\om_5$ and
$E_6\,\om_6$ --- with no disagreement: the dimension vectors agree in every
case, and the socle layers agree with the layers of (ii) after reversing the
degree, which is the expected shift between a projective generated in degree
$0$ and an injective with socle in degree $0$.  The validation set was chosen
to include vertices \emph{not} fixed by $\nu$ ($A_3$, $D_5\,\om_4$,
$E_6\,\om_5$, $E_6\,\om_6$), since at a $\nu$-fixed vertex the distinction
between $\Pia e_\ka$ and $\Pia e_{\nu(\ka)}$ is invisible and an error in the
Nakayama bookkeeping would not show.

Method (i) alone was then run exhaustively.  We note precisely what
``exhaustively'' covers:

\begin{itemize}
\item $E_6$: all six fundamental $\la$, all $1272$ nontrivial extremal
weights (orbit sizes $27,72,216,720,216,27$).  Complete.
\item $E_7$: all seven fundamental $\la$, all $17635$ nontrivial extremal
weights (orbit sizes $126,576,2016,10080,4032,756,56$).  Complete.
\item $E_8$: $\la=\om_1$ and $\la=\om_8$ only (orbit sizes $2160$ and $240$).
Both reach $\max_pc_p=2$ and contain no element with a $c$-trace entry $3$ at
all.  The remaining six fundamental weights of $E_8$ have orbits too large for
the present implementation and are \emph{not} covered.
\end{itemize}

\section{The nine examples}

\begin{convention}
Vertices are numbered as in Bourbaki: $E_n$ has the chain
$1-3-4-5-\cdots-n$ with vertex $2$ attached to vertex $4$, so vertex $4$ is
trivalent.  A weight $\mu$ is written in the basis of fundamental weights,
$\mu=(\langle\al_1^\vee,\mu\rangle,\dots,\langle\al_n^\vee,\mu\rangle)$.
Dimension vectors $v$ are written in the basis of simple roots, so
$\la-w\la=\sum_iv_i\al_i$.  A reduced word $s_{i_1}\cdots s_{i_\ell}$ acts with
$s_{i_\ell}$ first, and the $c$-trace is listed in that order, so $c_1$ is the
entry of the rightmost letter.  The displayed $V$ lists the socle layers of
$\soc_\sigma(I_\ka)$ from the bottom up, each layer as the multiset of its
colours.
\end{convention}

\begin{theorem}[computational]\label{thm:nine}
Let $\Gamma\in\{E_6,E_7\}$, let $\la=\om_\ka$ be any fundamental weight and
let $w$ be of minimal length in $wW_\la$.  Suppose the $c$-trace of $w$
attains the value $3$.  Then $T$ acts on $V$ with simple spectrum if and only
if $(\Gamma,\la,w)$ is one of the nine triples of Table \ref{tab:nine}.  Of
the $7094$ such triples in $E_6\cup E_7$, these nine are the only ones.
\end{theorem}

\begin{table}[ht]
\small
\begin{tabular}{@{}rllrl@{}}
\toprule
 & $\Gamma$, $\la$ & $w$ & $\ell(w)$ & $c$-trace \\
\midrule
(1) & $E_6$, $\om_4$ & $s_2s_4s_5s_6s_3s_4s_5s_1s_3s_4$ & $10$ & $(1^9,3)$ \\
(2) & $E_7$, $\om_3$ & $s_3s_4s_5s_6s_7s_2s_4s_5s_6s_1s_3s_4s_5s_2s_4s_3$ & $16$ & $(1^6,2,1^8,3)$ \\
(3) & $E_7$, $\om_3$ & $s_1s_3s_4s_5s_6s_7s_2s_4s_5s_6s_3s_4s_5s_2s_4s_3$ & $16$ & $(1^{15},3)$ \\
(4) & $E_7$, $\om_3$ & $s_1s_3s_4s_5s_6s_7s_2s_4s_5s_6s_1s_3s_4s_5s_2s_4s_3$ & $17$ & $(1^6,2,1^8,3,1)$ \\
(5) & $E_7$, $\om_4$ & $s_2s_4s_5s_6s_3s_4s_5s_1s_3s_4$ & $10$ & $(1^9,3)$ \\
(6) & $E_7$, $\om_4$ & $s_7s_2s_4s_5s_6s_3s_4s_5s_1s_3s_4$ & $11$ & $(1^9,3,1)$ \\
(7) & $E_7$, $\om_4$ & $s_6s_7s_2s_4s_5s_6s_3s_4s_5s_1s_3s_4$ & $12$ & $(1^9,3,1^2)$ \\
(8) & $E_7$, $\om_4$ & $s_5s_6s_7s_2s_4s_5s_6s_3s_4s_5s_1s_3s_4$ & $13$ & $(1^9,3,1^3)$ \\
(9) & $E_7$, $\om_5$ & $s_2s_4s_5s_6s_7s_3s_4s_5s_6s_1s_3s_4s_5$ & $13$ & $(1^{12},3)$ \\
\midrule
 & $\mu=w\la$ & $v=\la-w\la$ & $|v|$ & $V$, bottom layer first \\
\midrule
(1) & $(0,-3,0,2,0,0)$ & $(1,3,2,3,2,1)$ & $12$ & $\{2\}\,\{4\}\,\{2,3,5\}\,\{1,4,6\}\,\{2,3,5\}\,\{4\}$ \\
(2) & $(1,0,-3,2,0,0,0)$ & $(2,2,5,4,3,2,1)$ & $19$ & $\{3\}\,\{4\}\,\{2,3,5\}\,\{1,4,6\}\,\{3,5,7\}\,\{4,6\}\,\{2,3,5\}\,\{1,4\}\,\{3\}$ \\
(3) & $(-3,0,2,0,0,0,0)$ & $(3,2,3,4,3,2,1)$ & $18$ & $\{1\}\,\{3\}\,\{4\}\,\{2,5\}\,\{1,4,6\}\,\{3,5,7\}\,\{4,6\}\,\{2,5\}\,\{1,4\}\,\{3\}$ \\
(4) & $(-1,0,-2,2,0,0,0)$ & $(3,2,5,4,3,2,1)$ & $20$ & $\{1\}\,\{3\}\,\{4\}\,\{2,3,5\}\,\{1,4,6\}\,\{3,5,7\}\,\{4,6\}\,\{2,3,5\}\,\{1,4\}\,\{3\}$ \\
(5) & $(0,-3,0,2,0,0,1)$ & $(1,3,2,3,2,1,0)$ & $12$ & $\{2\}\,\{4\}\,\{2,3,5\}\,\{1,4,6\}\,\{2,3,5\}\,\{4\}$ \\
(6) & $(0,-3,0,2,0,1,-1)$ & $(1,3,2,3,2,1,1)$ & $13$ & $\{2\}\,\{4\}\,\{2,3,5,7\}\,\{1,4,6\}\,\{2,3,5\}\,\{4\}$ \\
(7) & $(0,-3,0,2,1,-1,0)$ & $(1,3,2,3,2,2,1)$ & $14$ & $\{2\}\,\{4,6\}\,\{2,3,5,7\}\,\{1,4,6\}\,\{2,3,5\}\,\{4\}$ \\
(8) & $(0,-3,0,3,-1,0,0)$ & $(1,3,2,3,3,2,1)$ & $15$ & $\{2,5\}\,\{4,6\}\,\{2,3,5,7\}\,\{1,4,6\}\,\{2,3,5\}\,\{4\}$ \\
(9) & $(0,-3,0,2,0,0,0)$ & $(1,3,2,3,3,2,1)$ & $15$ & $\{2\}\,\{4\}\,\{2,3,5\}\,\{1,4,6\}\,\{2,3,5,7\}\,\{4,6\}\,\{5\}$ \\
\bottomrule
\end{tabular}
\medskip
\caption{The nine extremal weights in $E_6\cup E_7$ whose $c$-trace attains
$3$ and whose $T$-fixed representation has simple spectrum.  All nine have
exactly one entry $c_p=3$.  Numbers (2) and (4) have in addition a single
entry $2$; the other seven are minuscule away from the single $3$.  Number
(4) is the only one that is not fully commutative.}
\label{tab:nine}
\end{table}

\begin{remark}[what the table does and does not pattern]\label{rem:pattern}
Three features are worth recording, and one tempting guess is false.

\emph{(a) The $3$ sits at or near the end, but ``at the end'' is not the
criterion.}  Four of the nine have $c$-trace exactly $(1,\dots,1,3)$.  It is
tempting to conjecture that among the elements with $\max_pc_p=3$, simple
spectrum holds precisely for those with $c$-trace $(1,\dots,1,3)$.  This is
false: of the nine, five have another shape, and conversely the condition is
not sufficient in the $\max_pc_p=2$ range, where $15$ elements of $E_6\cup
E_7$ have $c$-trace $(1,\dots,1,2)$ and fail simple spectrum.  We flag this
because the guess survives any inspection of minimal-length representatives
--- in each of the three families above, the shortest element does have
$c$-trace $(1,\dots,1,3)$ --- and is refuted only by the full list.

\emph{(b) $\la$ is always adjacent to or at the trivalent node.}  All nine
occur at $\la\in\{\om_3,\om_4,\om_5\}$, and in $E_6$ only at $\om_4$.  This is
consistent with the $E_8$ data: $\om_1$ and $\om_8$, the two fundamental
weights farthest from vertex $4$, admit no $c$-trace entry $3$ whatsoever.

\emph{(c) Items (5)--(8) and (9) extend (1) by minuscule letters.}  Example
(1) is the unique one in $E_6$; (5) is its image in $E_7$, with the same $w$,
the same $v$ in the first six coordinates and the same $V$; and (6), (7), (8)
are obtained from (5) by prepending $s_7$, $s_6s_7$, $s_5s_6s_7$, each step
appending a single minuscule letter and adding one box to $V$.  So the nine
are not nine unrelated accidents but four orbits of one phenomenon: a single
non-minuscule step of weight $3$ at the trivalent node, surrounded by
minuscule steps.

\emph{(d) $c_p\ge4$.}  There are $7560$ extremal weights in $E_6\cup E_7$ with
$\max_pc_p\ge4$ --- all of them in $E_7$ at $\la=\om_4$, all with
$\max_pc_p=4$ exactly --- and not one has simple spectrum.  We state this as a
computation and not as a theorem; it is the natural next thing to prove, and a
proof would have to explain why $3$ is possible and $4$ is not.
\end{remark}

\section{Example (1) in full}

The smallest of the nine is small enough to be checked by hand, and we give it
in full so that a reader can verify Table \ref{tab:nine} independently of any
code.

Take $\Gamma=E_6$, $\la=\om_4$ and
$w=s_2s_4s_5s_6s_3s_4s_5s_1s_3s_4$, of length $10$.  Acting with the rightmost
letter first on $\la=(0,0,0,1,0,0)$, the trace is
\[
\begin{array}{r@{\quad}c@{\quad}r@{\quad}l}
 p & i_p & c_p & \mu^{(p)} \\
 1 & 4 & 1 & (0,1,1,-1,1,0) \\
 2 & 3 & 1 & (1,1,-1,0,1,0) \\
 3 & 1 & 1 & (-1,1,0,0,1,0) \\
 4 & 5 & 1 & (-1,1,0,1,-1,1) \\
 5 & 4 & 1 & (-1,2,1,-1,0,1) \\
 6 & 3 & 1 & (0,2,-1,0,0,1) \\
 7 & 6 & 1 & (0,2,-1,0,1,-1) \\
 8 & 5 & 1 & (0,2,-1,1,-1,0) \\
 9 & 4 & 1 & (0,3,0,-1,0,0) \\
10 & 2 & 3 & (0,-3,0,2,0,0)
\end{array}
\]
so $\mu=w\la=(0,-3,0,2,0,0)$ and the single non-minuscule step is the last
one, at the trivalent neighbour $2$, where $\langle\al_2^\vee,\mu^{(9)}\rangle=3$.
Hence
\[
   v\ =\ \la-w\la\ =\ \al_1+3\al_2+2\al_3+3\al_4+2\al_5+\al_6,
   \qquad |v|=12 .
\]
Since $\mu\in W\la$ we have $(\mu,\mu)=(\la,\la)$, so $\delta(v)=0$ and $v$ is
extremal; by the uniqueness statement of the companion work, $I_\ka$ has at
most one submodule of dimension vector $v$, and exhibiting one identifies it
with $V$.  The weighted sequence $\sigma$ is the word read in trace order with
the last letter tripled,
\[
   \sigma\ =\ (4,3,1,5,4,3,6,5,4,2,2,2),
\]
and the iterated socle $\soc_\sigma(I_4)$ has layers, from the bottom,
\[
   \{4\},\quad \{2,3,5\},\quad \{1,4,6\},\quad \{2,3,5\},\quad \{4\},\quad \{2\} .
\]
Summing the layers gives $v$, each layer is multiplicity-free, and each colour
appears in distinct layers; so the spectrum is simple.  Note that the last
socle step, of colour $2$ and multiplicity $3$, contributes the three
occurrences of colour $2$ in three \emph{different} layers --- the second,
fourth and sixth from the bottom --- which is exactly why a step of weight $3$ need
not destroy simple spectrum.  A step of weight $c$ contributing $c$ boxes to a
\emph{single} layer is the failure mode, and it is what happens at all of the
other $479$ elements of $E_6$ with a $c$-trace entry $3$.

\section{Reproducing the computation}

The scripts are in the repository accompanying this note.  With
\verb|rootsys.py|, \verb|proj2.py|, \verb|wordtrace.py|, \verb|socle.py| and
\verb|dfs3.py| in the working directory:

\begin{itemize}
\item \verb|python3 socle.py| --- validates the socle method against the
submodule search on the worked examples, including the $D_4$ element that is
fully commutative without simple spectrum.
\item \verb|python3 socval.py E6 5| --- full cross-validation of the two
methods over one orbit; likewise \verb|A3 1|, \verb|D5 4|, \verb|E6 6|,
\verb|E6 4|, \verb|D4 2|.
\item \verb|python3 escan.py E6| and \verb|python3 escan.py E7| --- the
exhaustive scans, about $3$ seconds and about $3$ minutes respectively.
\item \verb|python3 enine.py| --- regenerates Table \ref{tab:nine}.
\item \verb|python3 escan.py E8 1 1| and \verb|python3 escan.py E8 8 8| ---
the two $E_8$ weights that are within reach.
\end{itemize}

All computations are linear algebra over $\mathbb{F}_p$, with $p$ set by the
environment variable \verb|DKKP| (default $7$).  The socle filtration is a
sequence of kernels of integer matrices, so a $p$-dependent answer would
require a $p$-torsion coincidence in one of those kernels.  We ran the
complete $E_6\cup E_7$ scan at $p=7$, $p=11$ and $p=32003$; all four counts of
\S\ref{sec:scope} and Remark \ref{rem:pattern}, and the list of nine itself,
were identical at all three.  This is evidence of $p$-independence, not a
proof of it.

\section{Open ends}

\begin{enumerate}[label=(\arabic*)]
\item $E_8$ at $\la=\om_2,\dots,\om_7$ is not covered.  The obstruction is the
orbit enumeration, not the socle computation: the latter costs about
$10$--$100$\,ms per element, so $\om_7$ (orbit $6720$), $\om_2$ (orbit
$17280$) and $\om_6$ (orbit $60480$) are a matter of minutes to hours, while
$\om_3,\om_4,\om_5$ run to the millions and need the search restricted by the
$c$-trace rather than run over the whole orbit.
\item Remark \ref{rem:pattern}(d) --- that $c_p\ge4$ never admits simple
spectrum --- is unproved, and is the sharpest statement the data suggests.
\item A conceptual description of the nine, replacing the list.  Remark
\ref{rem:pattern}(c) reduces them to four, and the mechanism isolated at the
end of \S3 --- whether the $c$ boxes of a weight-$c$ step land in $c$ distinct
socle layers --- is presumably the right invariant to make precise.
\item Whether the nine examples survive in characteristic $2$, where the
strand decomposition of the companion work is unavailable.
\end{enumerate}

\end{document}
